% ------------------------------------------------------------------------------------
%                               SEGNALI
% ------------------------------------------------------------------------------------

\chapter{Segnali}

Nelle reti di telecomunicazione, le informazioni sono \textit{trasmesse attraverso segnali}, 
ovvero dei \textit{modelli matematici di grandezze fisiche}.

\begin{definition}[SEGNALE]
    Un segnale è una \textbf{funzione} che associa un \textbf{valore} ad
    ogni \textbf{istante di tempo}:
    \[
        s(t) : \mathbb{R} \to \mathbb{C}
    \] 
\end{definition}

% ----- NOTA SULLA RAPPRESENTAZIONE -----

\section{Note sulla rappresentazione complessa}

In generale, si utilizza una \textbf{funzione complessa di variabile reale} per 
una questione di "comodità" matematica.
Il perchè appararirà più chiaro in futuro. Per ora ricordiamo l'uso 
della seguente notazione (usiamo $j$ per l'unità immaginaria):

\begin{itemize}
    \item Parte reale: $\text{Re}[s(t)]=s_R(t)$
    \item Parte immaginaria: $\text{Im}[s(t)] = s_I(t)$
    \item Complesso coniugato: $s^*(t)=\bar{s}(t) = s_R(t) - js_I(t)$
    \item Modulo: $|s(t)| = \sqrt{s_R^2(t) + s_I^2(t)} = \sqrt{s(t)s^*(t)}$
    \item Fase: $\text{arg}[s(t)] = \angle s(t) = \arctan(\frac{s_I(t)}{s_R(t)})$
\end{itemize}

% ----- TIPOLOGIE SEGNALI -----

\section{Tipologie di segnali}

Possiamo distinguere i segnali in base a:
\begin{itemize}
    \item \textbf{VALORI}
        \begin{itemize}
            \item \textbf{DISCRETI}: $s_R(t),s_I(t)\in \mathbb{Z}$
            \item \textbf{CONTINUI} : $s_R(t),s_I(t)\in \mathbb{R}$
        \end{itemize}
    \item \textbf{TEMPO}
        \begin{itemize}
            \item \textbf{DISCRETO}: $t\in \mathbb{Z}$
            \item \textbf{CONTINUO} : $t\in \mathbb{R}$
        \end{itemize}
\end{itemize}

Tale suddivisione è riassunta nella tabella~\ref{tb:01_01_tipologie_segnali}.

\begin{table}[!ht]
    \centering
    \renewcommand{\arraystretch}{1.6} 
    \begin{tabular}{|l|c|c|}
        \hline
        & Tempo discreto & Tempo continuo \\
        \hline
        
        Valori discreti & \cellcolor{orange!40}DIGITALI & no nome \\
        \hline
        
        Valori continui & no nome & \cellcolor{orange!40}ANALOGICI \\
        \hline
    \end{tabular}
    \caption{Tipologie di Segnali}
    \label{tb:01_01_tipologie_segnali}
\end{table}

La figura~\ref{fig:02_01_esempi_segnali} mostra un esempio di segnale analogico e di uno digitale.

\begin{figure}[!ht]
    \centering
    % --- Primo Plot: Segnale Analogico ---
    \begin{minipage}{0.48\textwidth}
        \centering
        \begin{tikzpicture}
            \begin{axis}[
                width=\textwidth,
                height=5cm,
                axis lines=middle,
                xlabel={$t$},
                ylabel={$s(t)$},
                xmin=0, xmax=2.2,
                ymin=-1.5, ymax=1.5,
                xtick={0, 0.5, 1, 1.5, 2},
                ytick={-1, 0, 1},
                grid=major,
                grid style={dashed, gray!30},
                title={Segnale Analogico},
                title style={font=\bfseries\small},
                xlabel style={below right},
                ylabel style={above left},
            ]
                \addplot[thick, nord10, domain=0:2, samples=200] {sin(deg(2*pi*x))};
            \end{axis}
        \end{tikzpicture}
    \end{minipage}\hfill
    % --- Secondo Plot: Segnale Digitale ---
    \begin{minipage}{0.48\textwidth}
        \centering
        \begin{tikzpicture}
            \begin{axis}[
                width=\textwidth,
                height=5cm,
                axis lines=middle,
                xlabel={$n$},
                ylabel={$s[n]$},
                xmin=-0.5, xmax=7.5,
                ymin=-0.5, ymax=1.5,
                xtick={0,1,2,3,4,5,6,7},
                ytick={0,1},
                grid=major,
                grid style={dashed, gray!30},
                title={Segnale Digitale},
                title style={font=\bfseries\small},
                xlabel style={below right},
                ylabel style={above left},
            ]
                \addplot[
                    ycomb, 
                    thick, 
                    nord11, 
                    mark=*, 
                    mark options={fill=nord11, scale=1.2}
                ] coordinates {
                    (0,1) (1,0) (2,1) (3,0) (4,0) (5,1) (6,1) (7,0)
                };
            \end{axis}
        \end{tikzpicture}
    \end{minipage}
    \caption{Un segnale analogico e un segnale digitale.}
    \label{fig:02_01_esempi_segnali}
\end{figure}

In particolare, per i segnali a \textbf{tempo discreto} definiamo il:

\begin{definition}[PERIODO DI SEGNALE]
    Il periodo di segnale $T_s$ è il periodo con il quale il segnale viene campionato.
    In tal senso risulta comodo usare la notazione seguente:
    \[
            t = nT_s, n\in \mathbb{Z} \implies s[n]\coloneqq s(nT_s)\in \mathbb{C}
    \] 
\end{definition}

% ----- ESEMPI SEGNALI -----

\section{Esempi di segnali}
\subsection{Impulso Rettangolare}

L'impulso rettangolare è così definito:

\begin{definition}[IMPULSO RETTANGOLARE]
    \[
        \text{rect}(t) \coloneqq \begin{cases}
            1 & |t| \le \frac{1}{2} \\
            0 & \text{altrimenti}
        \end{cases} 
    \] 
\end{definition}

\begin{figure}[!ht]
    \centering
    \begin{tikzpicture}[scale = 0.8]
        \begin{axis}[
            width=0.8\textwidth,
            height=5cm,
            axis lines=middle,
            xlabel={\(t\)},
            ylabel={\(s(t)\)},
            xmin=-1.1, xmax=1.1,
            ymin=-0.2, ymax=1.5,           
            xtick={-1, -0.5, 0, 0.5, 1},
            ytick={0, 1},
            grid=major,
            grid style={dashed, nord3!30},       
            axis line style={nord0},             
            tick style={nord0},                  
            tick label style={color=nord0},      
            xlabel style={below right, color=nord0},
            ylabel style={above left, color=nord0}
        ]
        \addplot[thick, nord10] coordinates {
            (-0.5, 0) (-0.5, 1) (0.5, 1) (0.5, 0)
        };
        \end{axis}
    \end{tikzpicture}
    \caption{Impulso rettangolare.}
\end{figure}

\subsection{Triangolo}

La funzione triangolo è definita come segue:

\begin{definition}[TRIANGOLO]
    \[
        \text{triang}(t) \coloneqq \begin{cases}
            1-|t| & |t| \le 1 \\
            0 & \text{altrimenti}
        \end{cases} 
    \] 
\end{definition}

\begin{figure}[!ht]
    \centering
    \begin{tikzpicture}[scale = 0.8]
        \begin{axis}[
            width=0.75\textwidth,
            height=5cm,
            axis lines=middle,
            xlabel={$t$},
            ylabel={$s(t)$},
            xmin=-2.5, xmax=2.5,
            ymin=-0.5, ymax=1.5,
            xtick={-2,-1,0,1,2},
            ytick={0,1},
            grid=major,
            grid style={dashed, gray!30},
            xlabel style={below right},
            ylabel style={above left},
        ]
            \addplot[thick, nord10] coordinates {
                (-2.5,0) (-1,0) (0,1) (1,0) (2.5,0)
            };
            
        \end{axis}
    \end{tikzpicture}
    \caption{Segnale triangolo.}
    \label{fig:02_03segnale_tri}
\end{figure}

\subsection{Esponenziale complesso}

Un esponenziale complesso è un segnale del tipo:

\begin{definition}[ESPONENZIALE COMPLESSO]
    \[
        s(t) = e^{j 2\pi f_{0}t} = \cos (2\pi f_{0}t) + j \sin (2\pi f_{0}t)
    \] 
\end{definition}

Notiamo che:

\begin{remark}
    E' un segnale \textbf{periodico}:
    \begin{itemize}
        \item Di \textit{frequenza} $f_{0}$
        \item Di periodo $T_{0} = \frac{1}{f_{0}}$
        \item Di pulsazione $\omega_{0} = 2\pi f_{0}$
        \item Di modulo $|e^{j 2\pi f_{0}t}|=1$
        \item Di fase $2\pi f_{0}t$
    \end{itemize}
\end{remark}

Un esempio di questo segnale è mostrato nella figura~\ref{fig:02_04:esponenziale_complesso}

\begin{figure}[!ht]
    \centering
    \begin{tikzpicture}[scale = 0.8]
        \begin{axis}[
            width=0.8\textwidth,
            height=0.65\textwidth,
            view={60}{30}, 
            axis lines=middle,
            xlabel={$t$},
            ylabel={$\mathrm{Re}$},
            zlabel={$\mathrm{Im}$},
            xmin=0, xmax=4*pi,
            ymin=-1.5, ymax=1.5,
            zmin=-1.5, zmax=1.5,
            xtick={0, pi, 2*pi, 3*pi, 4*pi},
            xticklabels={$0$, $\pi$, $2\pi$, $3\pi$, $4\pi$},
            ytick={-1,0,1},
            ztick={-1,0,1},
            grid=major,
            grid style={dashed, gray!30},
            xlabel style={below right},
            ylabel style={above left},
            zlabel style={above left},
        ]
            % Spirale principale (elica)
            \addplot3[
                domain=0:4*pi,
                samples=200,
                samples y=0, 
                variable=\t,
                thick,
                nord10,
            ] ({\t}, {cos(deg(\t))}, {sin(deg(\t))});
            
            % Proiezione sul piano t-Re 
            \addplot3[
                domain=0:4*pi,
                samples=200,
                samples y=0,
                variable=\t,
                dashed,
                nord14,
                opacity=0.5,
            ] ({\t}, {cos(deg(\t))}, {-1.5});
            
            % Proiezione sul piano t-Im 
            \addplot3[
                domain=0:4*pi,
                samples=200,
                samples y=0,
                variable=\t,
                dashed,
                nord11,
                opacity=0.5,
            ] ({\t}, {1.5}, {sin(deg(\t))});
            
            % Cerchio di proiezione sul piano Re-Im a t=0 
            \addplot3[
                domain=0:360,
                samples=100,
                samples y=0,
                variable=\a,
                dashed,
                nord3,
                opacity=0.4,
            ] ({0}, {cos(\a)}, {sin(\a)});
        \end{axis}
    \end{tikzpicture}
    \caption{Segnale esponenziale complesso}
    \label{fig:02_04:esponenziale_complesso}
\end{figure}

\subsection{Funzione sinc}

Un altro segnale molto importante è il seguente:

\begin{definition}[SINC]
    \[
        \text{sinc}(t) = \frac{\sin (\pi t)}{\pi t}
    \] 
\end{definition}

\begin{figure}[!ht]
    \centering
    \begin{tikzpicture}[scale = 0.8]
        \begin{axis}[
            width=0.85\textwidth,
            height=6cm,
            axis lines=middle,
            xlabel={$t$},
            ylabel={$\mathrm{sinc}(t)$},
            xmin=-5, xmax=5,
            ymin=-0.4, ymax=1.2,
            xtick={-5,-4,-3,-2,-1,0,1,2,3,4,5},
            ytick={-0.2,0,0.2,0.4,0.6,0.8,1},
            grid=major,
            grid style={dashed, gray!30},
            xlabel style={below right},
            ylabel style={above left},
            declare function={
                sinc(\x) = \x == 0 ? 1 : sin(deg(pi*\x))/(pi*\x);
            },
        ]
            % Plot della funzione sinc
            \addplot[
                thick,
                nord10,
                domain=-5:5,
                samples=400,
                smooth,
            ] {sinc(x)};
            
            % Evidenzia il punto in t=0 con un marker
            \addplot[
                only marks,
                mark=*,
                mark size=1.8pt,
                nord11,
            ] coordinates {(0,1)};
            
        \end{axis}
    \end{tikzpicture}
    \caption{Segnale sinc}
    \label{fig:02_05_segnale_sinc}
\end{figure}

Spesso ci si riferisce a questa funzione come "\textbf{sinc normalizzata}".

E' inoltre interessante notare che:

\begin{remark}
    Note sulla funzione sinc normalizzata:
    \begin{enumerate}
        \item $\lim_{t\to 0} \text{sinc}(t) = 1$
        \item La funzione si \textbf{annulla} per tutti gli 
            \textbf{interi non nulli}: $\text{sinc}(t)=0\forall k\in \mathbb{Z}\setminus \{0\}$
    \end{enumerate}
\end{remark}

% ----- OPERAZIONI SU SEGNALI -----

\section{Operazioni sui segnali}

\subsection{Ritardo/anticipo di un segnale}

Dato un segnale $s(t)$, possiamo \textit{traslarlo nel tempo} come segue:

\begin{definition}[TRASLAZIONE SEGNALE]
    \[
        \tilde{s}(t) \coloneqq s(t-t_{0})
    \] 
\end{definition}

In particolare, abbiamo:
\begin{itemize}
    \item \textbf{RITARDO} se $t_{0}>0$.
    \item \textbf{ANTICIPO} se $t_{0}<0$.
\end{itemize}

\begin{figure}[!ht]
    \centering
    \begin{tikzpicture}
        \begin{axis}[
            width=0.8\textwidth,
            height=5cm,
            axis lines=middle,
            xlabel={\(t\)},
            ylabel={\(s(t)\)},
            xmin=-1.5, xmax=3.5,
            ymin=-0.2, ymax=1.8, % Spazio sufficiente per le etichette
            xtick={-0.5, 0, 0.5, 1, 1.5, 2, 2.5},
            ytick={0, 1},
            grid=major,
            grid style={dashed, nord3!30},
            axis line style={nord0},
            tick style={nord0},
            tick label style={color=nord0},
            xlabel style={below right, color=nord0},
            ylabel style={above left, color=nord0}
        ]
        
        \addplot[thick, dashed, nord10] coordinates {
            (-0.5, 0) (-0.5, 1) (0.5, 1) (0.5, 0)
        };
        \node[nord10, anchor=south] at (axis cs: -0.5, 1.05) {\(s(t)\)};

        \addplot[thick, nord11] coordinates {
            (1.5, 0) (1.5, 1) (2.5, 1) (2.5, 0)
        };
        \node[nord11, anchor=south] at (axis cs: 2, 1.05) {\(\tilde{s}(t) = s(t-2)\)};

        \draw[<->, >=stealth, thick, nord3] (axis cs: 0, 0.6) -- (axis cs: 2, 0.6);
        \node[nord3, anchor=south] at (axis cs: 1, 0.6) {\(t_0 = 2\)};

        \end{axis}
    \end{tikzpicture}
    \caption{Traslazione del rettangolo.}
\end{figure}

\subsection{Scalamento di un segnale}

Nella trasmissione di un segnale, il ricevente ha spesso bisogno di un
\textbf{amplificatore} che \textit{scali il segnale}. Questo perchè
il segnale ricevuto risulta smorzato/attenuate durante la trasmissione.

\begin{definition}[SCALAMENTO SEGNALE]
    \[
        \tilde{s}(t) \coloneqq As(t)
    \] 
\end{definition}

Per un amplificatore, chiaramente vale $A>1$.

\begin{figure}[!ht]
    \centering
    \begin{tikzpicture}
        \begin{axis}[
            width=0.8\textwidth,
            height=5cm,
            axis lines=middle,
            xlabel={\(t\)},
            ylabel={\(s(t), \tilde{s}(t)\)},
            xmin=-1.1, xmax=1.1,
            ymin=-0.2, ymax=2.5, % Spazio per l'ampiezza 2
            xtick={-1, -0.5, 0, 0.5, 1},
            ytick={0, 1, 2},
            grid=major,
            % Applicazione della palette Nord
            grid style={dashed, nord3!30},
            axis line style={nord0},
            tick style={nord0},
            tick label style={color=nord0},
            xlabel style={below right, color=nord0},
            ylabel style={above left, color=nord0}
            ]

            % Segnale originale s(t) in blu tratteggiato
            \addplot[thick, dashed, nord10] coordinates {
                    (-0.5, 0) (-0.5, 1) (0.5, 1) (0.5, 0)
                };
            \node[nord10, anchor=south] at (axis cs: -0.2, 1.05) {\(s(t)\)};

            % Segnale scalato \tilde{s}(t) = 2s(t) in rosso continuo
            \addplot[thick, nord11] coordinates {
                    (-0.5, 0) (-0.5, 2) (0.5, 2) (0.5, 0)
                };
            \node[nord11, anchor=south] at (axis cs: 0.3, 2.05) {\(\tilde{s}(t) = 2s(t)\)};
        \end{axis}
    \end{tikzpicture}
    \caption{Scalamento dell'ampiezza di un segnale.}
\end{figure}

\begin{remark}
    In generale, $A$ può assumere qualsiasi valore reale,
    \textit{anche negativo} (in tal caso il segnale è riflesso
    rispetto all'asse delle ascisse)
\end{remark}

\subsection{Scalamento nel tempo}

Dato un segnale, oltre che scalarne l'ampiezza, possiamo scalarlo
anche nel tempo:

\begin{definition}[SCALEMENTO NEL TMEPO]
    \[
        \tilde{s}(t) \coloneqq s\left( \frac{t}{B} \right)
    \] 
\end{definition}

\begin{figure}[!ht]
    \centering
    \begin{tikzpicture}
        \begin{axis}[
            width=0.8\textwidth,
            height=5cm,
            axis lines=middle,
            xlabel={\(t\)},
            ylabel={\(s(t), \tilde{s}(t)\)},
            xmin=-1.8, xmax=1.8,
            ymin=-0.2, ymax=1.8, % Spazio per le etichette
            xtick={-1, -0.5, 0, 0.5, 1},
            ytick={0, 1},
            grid=major,
            % Applicazione della palette Nord
            grid style={dashed, nord3!30},
            axis line style={nord0},
            tick style={nord0},
            tick label style={color=nord0},
            xlabel style={below right, color=nord0},
            ylabel style={above left, color=nord0}
        ]
        
        % Segnale originale s(t) in blu tratteggiato
        \addplot[thick, dashed, nord10] coordinates {
            (-0.5, 0) (-0.5, 1) (0.5, 1) (0.5, 0)
        };
        \node[nord10, anchor=south] at (axis cs: 0.65, 0.25) {\(s(t)\)};

        % Caso |B| > 1 (Allungato): \tilde{s}(t) = s(t/2) in rosso
        \addplot[thick, nord11] coordinates {
            (-1, 0) (-1, 1) (1, 1) (1, 0)
        };
        \node[nord11, anchor=south] at (axis cs: -1.4, 0.25) {\(\tilde{s}(t) = s(\frac{t}{2})\)};

        % Caso |B| < 1 (Contratto): \tilde{s}(t) = s(2t) in verde
        \addplot[thick, nord14] coordinates {
            (-0.25, 0) (-0.25, 1) (0.25, 1) (0.25, 0)
        };
        \node[nord14, anchor=south] at (axis cs: -0.4, 1.05) {\(\tilde{s}(t) = s(2t)\)};

        \end{axis}
    \end{tikzpicture}
    \caption{Scalamento nel tempo di un segnale.}
\end{figure}

Ragionando su $B$ notiamo che:

\begin{remark}
    Distinguiamo 3 casi:
    \begin{itemize}
        \item Se $B<0$, il segnale è \textit{riflesso rispetto alle ordinate}.
        \item Se $|B|>1$ il segnale è \textbf{allungato}.
        \item Se $|B|<1$ il segnale è \textbf{contratto}.
    \end{itemize}
\end{remark}

\subsection{Combinazione di operazioni}

In generale, possiamo combinare tutte le operazioni viste sopra
per ottenere un nuovo segnale:

\begin{definition}[COMBINAZIONE DI OPERAZIONI]
    \[
        \tilde{s}(t) = As \left( \frac{t-t_{0}}{B} \right)
    \] 
\end{definition}

Come procediamo in tali casi? Procediamo come mostrato sotto:
\[
    s(t) \mapsto s_{1}(t) = s(t-t_{0}) \mapsto s_{2}(t) = s\left( \frac{t-t_{0}}{B} \right) 
    \mapsto s_{3} = As\left( \frac{t-t_{0}}{B} \right) = \tilde{s}(t)
\]

% ----- GRANDEZZE ASSOCIATE -----

\section{Grandezze associate ad un segnale}

\subsection{Valor medio}

Per un segnale a \textit{tempo continuo} $s(t)$, il suo \textit{valor medio} nell'intervallo
$[t_{1},t_{2}]$ è definito come segue:

\begin{definition}[VALOR MEDIO - TC]
    \[
        m_s(t_{1},t_{2}) = \frac{1}{t_{2}-t_{1}} \int_{t_{1}}^{t_{2}} s(t)dt
    \] 
\end{definition}

Se invece il segnale è a \textit{tempo discreto} abbiamo, in
un intervallo $[n_{1},n_{2}]$:

\begin{definition}[VALOR MEDIO - TD]
    \[
        m_s(n_{1},n_{2}) = \frac{1}{n_{2}-n_{1}} \sum_{n=n_{1}}^{n_{2}} s(nT)
    \] 
\end{definition}

Possiamo estendere il concetto di valor medio a tutto l'asse reale:

\begin{prop}[VALOR MEDIO SU TUTTO L'ASSE]
    Valor medio su tutto l'asse a tempo continuo:
    \[
        m_s = \lim_{u\to\infty} \frac{1}{2u} \int_{-u}^{u} s(t)dt
    \]
    Valor medio su tutto l'asse a tempo discreto:
    \[
        m_s = \lim_{k\to\infty} \frac{1}{2k} \sum_{n=-k}^{k} s(nT)
    \] 
\end{prop}

E' facile intuire che:

\begin{remark}[VALOR MEDIO DI SEGNALI DI DURATA FINITA]
    Se un segnale ha \textbf{durata finita} (es. $\text{rect}$), l'area
    totale è finita e, di conseguenza, il \textbf{valor medio è nullo}.
\end{remark}

\subsection{Energia}

Possiamo definire l'\textbf{energia} associata ad un segnale come segue.
Per un segnale a tempo continuo:

\begin{definition}[ENERGIA DI UN SEGNALE - TC]
    \[
        E_s = \int_{-\infty}^{+\infty} s(t)^2 dt
    \] 
\end{definition}

Per un segnale a tempo discreto:

\begin{definition}[ENERGIA DI UN SEGNALE - TD]
    \[
        E_s = T\sum_{n=-\infty}^{+\infty} s(nT)^2
    \] 
\end{definition}

E' utile osservare che l'energia di un segnale è \textbf{invariante} rispetto
a \textbf{traslazioni temporali}.

Inoltre, essendo l'energia del segnale una quantità quadratica, bisogna
fare attenzione all'energia trasportata dalla \textbf{somma di due segnali}.
Consideriamo per esempio due segnali a tempo continuo $s_{1}(t)$ e $s_{2}(t)$ 
(il caso discreto è analogo) e la loro somma $s(t)=s_{1}(t)+s_{2}(t)$:
Cominciamo calcolando il modulo al quadrato del segnale:
\[
    \begin{aligned}
        |s_{1} + s_{2}|^2 &= (s_{1}+s_{2})(s_{1}^* + s_{2}^*) \\ 
                   &= |s_{1}|^2 + |s_{2}|^2 + s_{1}s_{2}^* + s_{2}s_{1}^* 
                   =  |s_{1}|^2 + |s_{2}|^2 + (s_{1}s_{2}^*) + (s_{1}s_{2}^*)^* \\ 
                   &= |s_{1}|^2 + |s_{2}|^2 + 2 \text{Re}(s_{1}s_{2}^*)
    \end{aligned}
\]
Sostituendo nella definizione otteniamo allora:

\begin{warning}[ENERGIA DELLA SOMMA DI DUE SEGNALI]
    \[
    \begin{aligned}
        E_s &= \int_{-\infty}^{\infty} \left[ |s_{1}|^2 + |s_{2}|^2 + 2 \text{Re}(s_{1}s_{2}^*) \right] dt \\ 
            &= E_{s_{1}} + E_{s_{2}} + 2 \text{Re} 
            \left( \int_{-\infty}^{+\infty} s_{1}(t)s_{2}^*(t) dt\right)
    \end{aligned}
    \] 
\end{warning}

Possiamo quindi notare che:

\begin{remark}
    Se i due segnali hanno \textbf{supporto temporale disgiunto}, allora l'energia
    risultante è semplicemente la somma delle due energie:
    \[
        D_{s_{1}} \cap D_{s_{2}} = \emptyset \implies s_{1}s_{2}^* = 0 \forall t 
        \implies E_s = E_{s_{1}} + E_{s_{2}}
    \] 
\end{remark}

\subsection{Potenza media}

Possiamo definire anche la potenza media di un segnale in certo intervallo.
Per un segnale a tempo continuo $s(t)$ nell'intervallo $[t_{1},t_{2}]$:

\begin{definition}[POTENZA MEDIA - TC]
    \[
        P_s = \frac{1}{t_{2}-t_{1}} \int_{t_{1}}^{t_{2}} s^2(t)dt
    \] 
\end{definition}

Per un segnale a tempo discreto $s(nT)$ nell'intervallo $[n_{1},n_{2}]$:

\begin{definition}[POTENZA MEDIA - TD]
    \[
        P_s = \frac{1}{n_{1}-n_{2}} \sum_{n=n_{1}}^{n_{2}} s^2(nT)
    \] 
\end{definition}

Chiaramente, come prima abbiamo fatto con il valor medio, possiamo estendere il concetto
di potenza media su tutto l'asse $(-\infty,+\infty)$ semplicemente applicando il limite.

